\documentclass{ermacrs2025C}

\usepackage{times}
\portuguese

\title{A Matemática das Redes Neurais Artificias:\\Uma Visão Geométrica}

\author{Davi Picanço dos Reis}{Universidade Federal do Oeste do
Pará}{davipicancodosreis@gmail.com}
\author{Helaine Cristina Moraes Furtado}{Universidade Federal do
Oeste do Pará}{helaine.furtado@ufopa.edu.br}

\abstract{As Redes Neurais Artificiais tornaram-se essenciais para a
  ciência e tecnologia nas últimas décadas, ganhando grande
  popularidade. Contudo, seu funcionamento técnico costuma ser
  restrito a especialistas, o que gera uma percepção "mágica" no público geral.
  Nesse contexto, este trabalho apresenta os fundamentos matemáticos e
  a interpretação geométrica dessas redes, visando promover uma
  compreensão intuitiva e introdutória sobre o tema.
}

\keyword{Redes Neurais Artificiais}
\keyword{Machine Learning}
\keyword{Inteligência Artificial}
\keyword{Álgebra Linear}

\begin{document}

\maketitle

\section{Discussão}

Uma Rede Neural Artificial (RNA) é um modelo computacional inspirado
pelo funcionamento do cérebro humano. É formada por vários neurônios
conectados que transferem informação entre si. E o seu principal
objetivo é desenvolver a capacidade de aprender com a experiência, ou
seja, aprender diretamente dos dados.

O Perceptron é o modelo mais simples de neurônio, ele é a unidade
computacional básica de uma RNA. Ele é responsável por receber
dados, na forma de vetor, fazer o processamento e retornar um valor
real. Assim, o seu processamento pode ser representado pela equação:
\begin{equation}
  y(\mathbf x)=\varphi(\mathbf w^T\cdot \mathbf x+b)
\end{equation}

O vetor $\mathbf x$ é o vetor com os dados de entrada, $\mathbf w$ é
o vetor de pesos, ou parâmetros, que também tem dimensão $n$, $b$ é o
\textit{bias}, $\varphi$ é chamada de função de ativação e $y(\mathbf
x)$ é o resultado do processamento ou a ativação do neurônio, dada
uma entrada de dados.

Agora, utilizando uma intuição geométrica, consideremos as seguintes
afirmações como verdadeiras:
\begin{itemize}
  \item O produto escalar entre dois vetores mede a "similaridade"
    entre dois vetores;
  \item Dados podem ser vistos como pontos em um espaço;
  \item Um vetor e um ponto determinam um hiperplano no espaço.
\end{itemize}

Por enquanto, tomemos a função de ativação $\varphi$ como a função
identidade, de modo que não afete o resultado do processamento.
Assim, considerando uma RNA já treinada, temos um vetor $\mathbf w$
já definido. Levando em consideração os pontos podemos ter até 2
níveis de uma interpretação geométrica.

Primeiramente consideremos que o vetor de pesos $\mathbf w$ já
treinado aponta para certa posição nesse espaço de dados, e quando
fazemos o produto escalar do vetor de entradas $\mathbf x$ com ele o
que estamos medindo é a similaridade entre eles, ou seja, de forma
mais intuitiva, se eles estão próximos da mesma posição. Portanto,
podemos dizer que um neurônio pode verificar, nessa representação de
dados no espaço, se os dados de entrada estão próximos ao padrão
esperado pelo neurônio.

A segunda abordagem leva em consideração uma equação do plano:
$p(x,y,z)=ax+bx+cx+d$. Podemos notar, ainda desconsiderando a função
de ativação, que a equação (1) também é uma equação de plano, mas em
um espaço multidimensional, ou seja, um hiperplano.
\begin{equation}
  y(x_1,x_2,...,x_n)=w_1x_1+w_2x_2+...+w_nx_n+b
\end{equation}

Nesse caso os $w_1,w_2,...,w_n,b$ seriam os coeficientes desse
hiperplano. Ou seja, um neurônio já treinado determina um hiperplano
no espaço de dados, e isso somado ao fato que os dados são pontos
nesse espaço podemos dizer que o hiperplano separa os dados que tem
certa semelhança dos que não tem.

Entretanto, a gama de problemas lineares, que podem ser separados com
uma reta, são muito menores que problemas não lineares. Assim, a
função proporciona uma forma de torna o processamento menor linear.
De forma intuitiva, o que a função de ativação faz é transformar,
dobrar, o espaço de modo que os dados se tornem linearmente
separáveis, que podem ser separados por uma linha reta.

Assim, temos a interpretação geométrica completa do processamento do
neurônio, um neurônio é responsável por separar o espaço de dados em
dois por meio de uma superfície, de modo a categorizar os dados que
contribuem ou não para o padrão esperado pelo neurônio.

Podemos ainda extender essa interpretação para uma RNA com mais
neurônios e camadas. Nesse caso cada neurônio fica responsável por
segmentar os dados de uma forma diferente, e criando no espaço de
dados vários subespaços que podem ser categorizados de acordo com o
esperado pela rede.

\printreferences

\end{document}
